Considere
\[
-\big[(1+x)y'\big]'=\lambda(1+x)y,\qquad0<x<2,\qquad y'(0)=y'(2)=0.
\]
Una autofunción debe ser no idénticamente nula. Determine si \(\lambda=0\) admite autofunciones y describa todas las que correspondan a ese valor, sin buscar el espectro completo. Explique qué libertad queda en su amplitud después de imponer las condiciones de contorno. Si existe un modo para \(\lambda=0\), seleccione el que sea positivo y tenga norma uno para
\(\|y\|_w^2=\int_0^2(1+x)y(x)^2\,dx\).
